\begin{abstract}
Can a useful sentiment classifier be extracted from a language model with
\emph{zero} human labels, a handful of generated stories, and three lines
of linear algebra? We show that the answer is yes. We introduce \ours, a
fully unsupervised recipe that generates nine short stories for each of nine
emotion categories (81 LLM-generated paragraphs, no polarity labels),
mean-pools the final-layer residual stream, and takes the first principal
component of the nine category centroids as a \emph{valence axis}.
Projecting held-out text onto this single 1\,536-dimensional direction
gives AUC \sstauc{} on SST-2, \imdbauc{} on IMDB, \yelpauc{} on Yelp
Polarity, and \twauc{} on Twitter SemEval-17, correlates
\huliur{} with the Hu~\& Liu 6{,}789-word lexicon and \warrinerr{}
with the Warriner continuous valence norms --- all from the \emph{same} 81
stories and no supervised training. Applying the \emph{English} axis
zero-shot to non-English held-out stories produces transfer AUC $0.949$
(Spanish), $0.991$ (French), $0.972$ (German), and $1.000$ (Chinese),
confirming that the direction is cross-lingual at extraction. A supervised
logistic regression on the same hidden features needs $N\!\approx\!100$
SST-2 labels to match our zero-label AUC and the axis is stable across
model scale (Qwen2.5-1.5B and 7B both give $|r|\!=\!0.494$ on EmoBank).
The method is robust in a second sense: it is the most data-efficient
one-dimensional summary of valence we can find. A full-rank logistic
regression on the residual stream reaches AUC $0.911$; ablating our axis
from the features changes that to $0.910$, showing that valence is
richly encoded across many directions and we merely extract the single
most useful one under a zero-label constraint. The same PCA recipe
\emph{also} outperforms the standard contrastive mean-difference axis of
Zou et al.\ and Arditi et al.\ ($0.868$ vs.\ $0.753$, $+0.115$ AUC) without
requiring the operator to pre-specify a positive/negative split. Layer-wise
probing reveals a sharp ``V-shape'': valence appears at layer 1
(embedding-level), fades through the context-integration middle, then
re-crystallizes at the final two transformer blocks. This pattern is
universal across Qwen, GPT-NeoX, and BLOOM families. We position \ours{}
as a simple, deployment-oriented probe: zero labels, nine languages
tested, one axis, seventeen evaluations, and \speedup{} faster than direct
prompting.
\end{abstract}
